\chapter{Integer structure of the constant readers}
\label{chap:canales-enteros-constantes}
\index{mathematical constants!integer channels}
\index{digital root!residue and carry}
\index{semicrown}

The preceding generations distinguish the function of each constant. This
block proves the integral structure of their readers. The quantities
\(729,271,315\) and \(161\) are constructed as \emph{integral channels relative
to declared typed rules}; they do not enter as decimal approximations. Their
statuses are not identical: \(729\) is the quadratic power of channel
\(27\); \(271\) and \(315\) receive a nonadic realization constructed from
the published TPK phase partition and a second realization relative to the
Steiner lift; and
\(161\) remains, within the present scope, an integral identity without a
set-theoretic realization of its own. The catalogue is then classified through
published predicates and gauges. Only at the end may an Archimedean chart
act to produce decimal expansions.

This order separates two distinct operations. Seeking the integer \(271\) because it recalls the first digits of \(e\) would be a retrospective reading. Finding a TPK region and calling it \(e\) without constructing the type of return it represents would be a designation without a map. The integer channels, the conditioned orbital classification, and the dictionary relating them are therefore kept separate. The isolated integer channel selects no constant; selection belongs to the enriched TPK reader, which incorporates the orbit, orientation, survival, and boundary memory without receiving as input the decimal value that it ultimately evaluates.

\section{Visible residue, digital root and crown memory}

For \(n>0\), the nonadic digital root is the section
\[
 \operatorname{dr}_9(n)
 =\begin{cases}
 9,&n\equiv0\pmod9,\\
 n\bmod9,&n\not\equiv0\pmod9.
 \end{cases}
\]
Reduction modulo nine conserves the visible class; the section by
\(\{1,\ldots,9\}\) also conserves the choice of the null representative.
We define the lifted quotient
\[
 q_9^\sharp(n):=\frac{n-\operatorname{dr}_9(n)}9
 \in\mathbb Z_{\geq0}.
\]
One obtains the exact decomposition, which is not the usual remainder convention
in \(\{0,\ldots,8\}\),
\[
 n=9q_9^\sharp(n)+\operatorname{dr}_9(n).
\]
This separates two readings: \(\operatorname{dr}_9(n)\) records the visible
value, and \(q_9^\sharp(n)\) records how many complete crowns precede the
chosen visible representative. This is the exact form of the colloquial
image according to which the modulus compacts and the carry preserves what
compactification would lose.

In APP the additive and multiplicative leaves produce
\[
 \sum\Sigma=405,
 \qquad
 \sum\Pi=459,
 \qquad
 \Delta:=\sum\Pi-\sum\Sigma=54.
\]
The sum of the additive quotients is
\[
 S:=\sum Q^+=45.
\]
We define
\[
 \kappa=\frac{\Delta}{2}=27,
 \qquad
 \Theta=5\kappa=135.
\]
All these numbers come from the orthogram and its two evaluations. No constant has been used yet.

The central block
\[
 B_c=\begin{pmatrix}7&2\\2&7\end{pmatrix}
\]
produce four ordered words:
\[
 72,\qquad27,\qquad77,\qquad22.
\]
Satisfy
\[
 72+27=77+22=99,
 \qquad
 72-27=45,
 \qquad
 77-22=55=54+1.
\]
The center does not remain isolated: its two complementary readings
project onto the boundary \(99\), while their differences recover the
additive crown \(45\) and the pointed defect \(54+1\).

\section{Theorem of the four integer channels}

\begin{theorem}[Typed identities of the four channels]
\label{thm:derivacion-canales-enteros}
Relative to the composition rules
\(\kappa\mapsto\kappa^2\), \(\Theta\mapsto2\Theta+1\),
\((\Theta,S)\mapsto2\Theta+S\), and
\((\Theta,\kappa)\mapsto\Theta+\kappa-1\), the four integer
channels are defined
\[
\begin{aligned}
 C_\alpha&=\kappa^2=729,\\
 C_e&=2\Theta+1=271,\\
 C_\pi&=2\Theta+S=315,\\
 C_\varphi&=\Theta+\kappa-1=161.
\end{aligned}
\]
Moreover,
\[
\begin{aligned}
 C_e&=5\Delta+1=10\cdot27+1
      =4\Delta+(\Delta+1),\\
 C_\pi&=5\Delta+S=7S,\\
 C_\varphi&=3\Delta-1
      =4\Delta-(\Delta+1),
\end{aligned}
\]
and the closures are satisfied
\[
 \boxed{1000=C_\alpha+C_e=729+271}
\]
and
\[
 \boxed{C_\pi+C_e+C_\varphi-3=744=729+15.}
\]
The \(\alpha\) channel also has the central reading
\(729=10\cdot72+9\), and the \(e\) channel has the reading
\(271=10\cdot27+1\).
\end{theorem}

\begin{proof}
Substituting \(\Delta=54\), \(S=45\), \(\kappa=27\), and
\(\Theta=135\) yields the four formulas. The alternative expressions are
integer identities. The first completion is obtained from
\(27^2+2(5\cdot27)+1=729+270+1=1000\); the second, from
\(315+271+161-3=744=729+15\). The two decimal divisions follow from
the central words \((72,27)\) and their remainders already fixed by the
cubic completion.
\end{proof}

The equality \(1000-3^6=271\) is universal for that cubic ambient space. By
itself it does not distinguish \(e\): in an ambient space of cardinality \(B\), constructed with
an alphabet of cardinality \(q\) in dimension \(d\), the formal complement is
\(B-q^d\). The additional information in this section is that \(271\) also admits
a realization as a pointed return and the central reading \(27|1\).
Its association with the orbit named \(e\) belongs to the subsequent typed
dictionary, not to the isolated subtraction.

\section{Semicorona nonadic relative to the TPK phase rule}

The minimal motion rule of the TPK emitter does not treat the nine instants
of a window as an interchangeable list. Relative to the published phase
origin, it defines
\[
 D_9=\{1,\ldots,9\},
 \qquad
 H_+=\{1,4,7\},
 \qquad
 H_-=\{2,5,8\},
 \qquad
 H_0=\{3,6,9\}.
\]
The instants of \(H_+\) activate the additive reading, those of \(H_-\) the
multiplicative reading, and those of \(H_0\) preserve the neutral channel. Set
\[
 H_{\rm act}=H_+\sqcup H_-,
 \qquad
 O_{\rm ph}=\{1,\ldots,8\},
 \qquad
 \star_{\rm ph}=9,
\]
and let \(P_2(H_{\rm act})\) be the set of unordered pairs of active
phases. We define
\[
\begin{aligned}
 \mathcal F^{\rm ph}_6&=H_{\rm act}\times P_2(H_{\rm act}),\\
 \mathcal F^{\rm ph}_8&=O_{\rm ph}\times P_2(H_{\rm act}),\\
 \Theta_{\rm ph}&=D_9\times P_2(H_{\rm act}).
\end{aligned}
\]

\begin{theorem}[Semicrown Relative to the TPK Phase]
\label{thm:semicorona-fase-tpk}
Relative to the published window origin and to the terminal return
\(9\to1\), it holds
\[
 |\mathcal F^{\rm ph}_6|=90,
 \qquad
 |\mathcal F^{\rm ph}_8|=120,
 \qquad
 |\Theta_{\rm ph}|=135.
\]
The oriented boundary, the pointed return, the neutral lateral boundary, and the
lateral closure are, respectively,
\[
\begin{aligned}
 \partial_{\rm ph}&=\Theta_{\rm ph}\times\{-1,+1\},\\
 E_{\rm ph}&=\partial_{\rm ph}\sqcup\{\infty\},\\
 L_{\rm ph}&=H_0\times P_2(H_{\rm act}),\\
 P_{\rm ph}&=\partial_{\rm ph}\sqcup L_{\rm ph},
\end{aligned}
\]
and satisfy
\[
 \boxed{|\partial_{\rm ph}|=270,\quad
 |E_{\rm ph}|=271,\quad |L_{\rm ph}|=45,\quad
 |P_{\rm ph}|=315.}
\]
The construction is covariant under a change of origin: translating the window
by \(C_9\) simultaneously translates all strata and changes neither their
cardinalities nor their incidences.
\end{theorem}

\begin{proof}
The sign rule gives six active phases and three neutral ones, so that

\[
 |P_2(H_{\rm act})|=\binom62=15.
\]

Multiplying by \(6\), \(8\), and \(9\) produces \(90\), \(120\), and \(135\).
The orientation doubles \(135\); the pointed return adds one element; and the
three neutral phases contribute \(3\cdot15=45\) lateral positions.
Covariance follows because a translation of \(D_9\) conjugates the phase rule
and transports each Cartesian product by a bijection.
\end{proof}

This theorem substantively strengthens the cardinal identity \(270+1\): the
order \(6\subset8\subset9\) is no longer chosen from among \(9!\) possible orders,
but instead arises from the TPK phase clock and its published return.
Canonicity is relative to the origin of the window; changing it does not alter the
object up to the natural action of \(C_9\).

\section{Semicorona relative to the hexad-octad elevation and pointed return}

Fix the dodecad, the alignment, and a hexad
\(H\in\mathcal H(D)\) of the construction of
\cref{sec:w12-w24-elevacion}. By
\cref{thm:elevacion-unica-hexada}, there exists a unique octad \(O_H\) such that
\(O_H\cap D=H\). Write
\[
 P_2(H)=\{\{a,b\}:a,b\in H,\ a\ne b\},
 \qquad |P_2(H)|=\binom62=15,
\]
and add to \(O_H\) a return state \(\star\):
\[
 \widehat O_H=O_H\sqcup\{\star\},
 \qquad |\widehat O_H|=9.
\]
We define
\[
 \Theta_H=\widehat O_H\times P_2(H).
\]
Within it appear the two combinatorial layers.
\[
 \mathcal F_6=H\times P_2(H),
 \qquad
 \mathcal F_8=O_H\times P_2(H).
\]

\begin{theorem}[Realization of the semicrown]
\label{thm:realizacion-semicorona}
The preceding objects satisfy
\[
 |\mathcal F_6|=90,
 \qquad
 |\mathcal F_8|=120,
 \qquad
 |\Theta_H|=135.
\]
The oriented boundary
\[
 \partial_H=\Theta_H\times\{-1,+1\}
\]
has \(270\) states. Adjoining one return point,
\[
 E_H=\partial_H\sqcup\{\infty\},
\]
has \(271\) states. The boundary lateral of the stratum of six points is
\[
 L_H=\Theta_H\setminus\mathcal F_6,
 \qquad |L_H|=45,
\]
and the lateral-oriented closure
\[
 P_H=\partial_H\sqcup L_H
\]
tiene \(315\) estados.
\end{theorem}

\begin{proof}
The cardinals are
\[
 6\binom62=90,
 \qquad
 8\binom62=120,
 \qquad
 9\binom62=135.
\]
The orientation doubles \(135\), adjoining one point adds a unit, and
the difference between \(\Theta_H\) and the six-point stratum contains
\((9-6)15=45\) elements. From this result \(270,271\) and
\(270+45=315\).
\end{proof}

Here the two integers have different set-theoretic types. \(271\) is the
cardinality of an oriented boundary to which one point has been adjoined; that
fact alone does not make it the orbit of \(e\). \(315\) is the cardinality of that
boundary together with the \(45\) lateral positions not belonging to the
six-point stratum. The inclusion \(H\subset O_H\) used here is indeed the
Steiner lift already proved, relative to the fixed dodecad and alignment. The
semicorona then adjoins the point \(\star\), the orientation, and the lateral
boundary. This construction does not yet provide the nonadic order induced by
TPK or an intertwiner between the faces of the cube and Steiner incidences;
these are additional maps.

The two semicrowns share the same cardinal profile, but they must not
therefore be identified. An unmarked isomorphism between the chains
\[
 H_{\rm act}\subset O_{\rm ph}\subset D_9
 \quad\text{and}\quad
 H\subset O_H\subset\widehat O_H
\]
that preserves the return point may freely permute the six elements of the first stratum and the two elements added in the second. There exist
\[
 \boxed{6!\,2!=1440}
\]
chain gauges. The phase--Steiner intertwiner must select one of
them by means of additional incidences; equality of cardinalities does not
construct it.

\section{Dihedral action and classification under declared predicates}

For an emission
\(U=(u_1,\ldots,u_6)\in\mathcal U_6\), we define the visible word
\[
 \Psi(U)=((-u_1)\bmod3,\ldots,(-u_6)\bmod3)
 \in\mathbb F_3^6.
\]
Let the coordinate permutations be explicitly defined by
\[
 r\cdot(u_1,u_2,u_3,u_4,u_5,u_6)
 =(u_3,u_1,u_2,u_5,u_6,u_4),
\]
\[
 s\cdot(u_1,u_2,u_3,u_4,u_5,u_6)
 =(u_4,u_5,u_6,u_1,u_2,u_3).
\]
Generan
a grupo diedral \(D_3\simeq S_3\). The catalogo verifica
\[
 r\mathcal U_6=\mathcal U_6,
 \qquad
 s\mathcal U_6=\mathcal U_6,
 \qquad
 \Psi(gU)=g\Psi(U),
\]
and preserves the microscopic multiplicities.

The action divides the \(243\) visible words into \(38\) orbits of size
six and five orbits of size three. On each word \(w\), write
\[
 n(w)=|\Psi^{-1}(w)|,
 \qquad
 M(w)=\sum_{\Psi(U)=w}|F^{-1}(U)|.
\]

We now fix three published predicates. Predicate \(P_\pi\) requires
\(n(w)=5\), \(M(w)=1008\), multiplicities
\((432,144,144,144,144)\), and tritic census \((2,3,1)\). Predicate
\(P_e\) requires, over the complete orbit, singleton emissions of
microscopic mass \(144\), union of values 
\(\bigcup_{w\in\mathcal O}\operatorname{diag}_c(w)=\{0,3,6\}\), and census
\((2,3,1)\). Predicate \(P_\varphi\) preserves the
first three conditions of \(P_e\) and requires balanced census \((2,2,2)\).
The orientation gauges are \(\operatorname{diag}_c=6\) and the cardinal
class is \(ES\).

\begin{theorem}[Finite orbital classification conditioned]
\label{thm:selectores-orbitales-constantes}
In the reconstructed catalog, the preceding predicates select the
following unique orbits without consulting digits of any constant.
\begin{enumerate}
  \item There is a unique orbit of six words for which, in every
word, \(n(w)=5\), \(M(w)=1008\), and the five multiplicities are
\(432,144,144,144,144\). All its words have trit census
\((n_0,n_1,n_2)=(2,3,1)\). It is the dihedral orbit of the regional type of
\(\pi\).
  \item Between the orbits radical of emissions one-point and mass
  microscopic \(144\) who union orbital of values
  \(\operatorname{diag}_c\) is \(\{0,3,6\}\), there exists to unique
  orbit with census \((2,3,1)\). Is the orbit of return of \(e\).
  \item In the same sector there is a unique balanced orbit with census
  \((2,2,2)\). It is the proportional orbit of \(\varphi\).
\end{enumerate}
The orientation declared by \(\operatorname{diag}_c=6\) and the cardinal class
\(ES\) select, respectively,
\[
 \boxed{
 w_\pi=010211,
 \qquad
 w_e=201101,
 \qquad
 w_\varphi=121200.
 }
\]
\end{theorem}

\begin{proof}
The exhaustive reconstruction of the \(468\) emissions first verifies
closure of the action and groups the \(243\) words into orbits. It then applies
the indicated predicates to all of them. The first predicate leaves a single
orbit; the two singleton radical-fiber predicates respectively leave an
orbit of census \((2,3,1)\) and one of census \((2,2,2)\). Finally, the
orientation gauge is applied to each orbit, leaving a unique word.
None of the selector predicates contains the decimal expansions of
\(\pi,e\) or \(\varphi\).
\end{proof}

\subsection{First intrinsic encoder and its limit}

In the stratum of visible words containing exactly two zeros, there is
an encoder that does arise from the TPK datum itself:
\[
 c_0(w)=\{j\in\{1,\ldots,6\}:w_j=0\}
 \in P_2(\{1,\ldots,6\}).
\]
It is \(D_3\)-equivariant because a coordinate permutation transports the zero support. On the orbits classified by the
\cref{thm:selectores-orbitales-constantes}, its images are
\[
\begin{aligned}
 c_0(\mathcal O_\pi)
 &=\{\{1,2\},\{1,3\},\{2,3\},\{4,5\},\{4,6\},\{5,6\}\},\\
 c_0(\mathcal O_e)
 &=\{\{1,6\},\{2,5\},\{3,4\}\},
\end{aligned}
\]
where, on the second line, each pair has two preimages within the orbit.
The orbit of \(\varphi\) shares the first support type with \(\pi\), but
is separated from it by its singleton fiber, its multiplicity, and its TRIT\@ census.

The map \(c_0\) is the first rung of the register--pair encoder; it is not the complete family \(c_k\). In deep blocks, (0,1,3) or (4) zeros occur, so the same formula does not take values in \(P_2\) at every depth. Furthermore, \(\{1,\ldots,6\}\) is here the hexad of coordinates of the word, not an already identified Witt hexad. Prolonging \(c_0\) requires a phase-dependent rule and the corresponding incidence intertwiner.

The canonical microfiber of \(\pi\) remains exactly
\[
 \mathscr R_\pi^{\rm micro}=\Psi^{-1}(010211),
\]
formed by five distinct regions \(U_6\). The other five words of the
dihedral orbit are gauge images of that orbital form; they do not constitute
five new canonical microfibers of \(\pi\). Analogously, the names
\(e\) and \(\varphi\) are applied after fixing the published predicates and gauge.
The enumeration is exhaustive \emph{relative} to those
premises. Predicates and gauge are declared components of the enriched TPK
reader: they are reconstructed not from the three coarse statistics
appearing below, but from the complete genealogy of the emission,
its orientation, its boundary, and its survival.

The sum of the six integers of an emission, its residue modulo nine, and the
number of even entries do not suffice for this theorem. The oriented emission
of \(e\) and three of the five regions of \(\pi\) share exactly
\[
 \sum_j u_j=3316,
 \qquad
 \sum_j u_j\equiv4\pmod9,
 \qquad
 \#\{j:u_j\text{ even}\}=4.
\]
The distinction requires the istructura of the fiber, the multiplicity, the sector
radical and the orbit complete. This ablation demonstrates by que the parity and
the root digital are ingredients of the reader, but not replace the
transport TPK\@.

\section{Typed correspondence \texorpdfstring{\(271\to e\) and
\(315\to\pi\)}{271 to e and 315 to pi}}

The
\cref{thm:semicorona-fase-tpk,thm:realizacion-semicorona,thm:selectores-orbitales-constantes}
are independent: the first constructs
the semicrown from the phase clock, the second provides its relative Steiner
realization, and the third classifies orbits under fixed predicates. To
compose the first and third results, we declare the
type dictionary that the reader must preserve:
\[
\begin{array}{c|c}
\text{internal type}&\text{orbital type}\\
\hline
\text{bidirectional pointed return}&
\text{single-point radical fiber with the census of the regional branch}\\
\text{boundary-plus-lateral closure}&
\text{closure orbit with five regions}.
\end{array}
\]

\begin{corollary}[Typed, orbit-relative associations]
\label{cor:lectores-enteros-relativos}
Relative to the dictionary, the predicates \(P_e,P_\pi\), and the
preceding orientation gauges, the orbital classification produces the
conditioned associations
\[
 \boxed{|E_{\rm ph}|=271\rightsquigarrow \mathcal O_e
        \longrightarrow w_e,}
 \qquad
 \boxed{|P_{\rm ph}|=315\rightsquigarrow \mathcal O_\pi
        \longrightarrow w_\pi.}
\]
In the second case, the oriented representative \(w_\pi=010211\) determines the
canonical five-region microfiber \(\Psi^{-1}(010211)\). The associations are obtained
before Archimedean evaluation and do not use the decimal values of the
constants.
\end{corollary}

\begin{proof}
By the orbital theorem, each row of the dictionary possesses a unique orbit
satisfying the declared predicate. Classification is first performed at the
level of the complete dihedral orbit; the oriented gauge then selects the
indicated representative. Uniqueness is therefore relative to the
dictionary, predicates, and gauges fixed.
\end{proof}

The channel \(161=\Theta+\kappa-1=3\Delta-1\) remains the integer identity.
Although \(P_\varphi\) classifies a balanced orbit and the published gauge
selects \(w_\varphi=121200\), this section does not construct a
set-theoretic object of cardinality \(161\) whose type induces that orbit. Therefore, no
reader \(161\to\varphi\) is asserted here.

The corollary depends on the type dictionary because cardinalities, taken in
isolation, do not define a function into the catalogue. In the complete reader,
that dictionary is realized through the nine monodromy relations of the
\cref{chap:extensiones-dinamica-cociente}, together with orientation and
regional survival. The important consequence is affirmative: once
the structural reader has been fixed, the evaluation consults no decimal
approximation of the constant.

\subsection{Joint dependence of the channels on the product selector}

The executable factorization of the catalog and the parametrization of its additive signature give
\[
 \mathcal U_6\simeq\mathcal S_+\times\mathcal S_\times,
 \qquad
 \mathcal S_+\simeq D_9\times\{\pm1\}.
\]
On the other hand,
\[
 \partial_{\rm ph}
 \simeq(D_9\times\{\pm1\})\times P_2(H_{\rm act}).
\]
Thus, a functional attempt at descent that preserves phase and orientation identically reduces to the map
\[
 q_\times:\mathcal S_\times\longrightarrow P_2(H_{\rm act}).
\]
This is the field absent from the catalogue: the published rows end at the
product signature \(E_{\rm sig}\) and do not contain a class
\(P_2(H_{\rm act})\).

\begin{proposition}[Obstruction to the isolated product selector]
\label{prop:nogo-selector-producto-d3}
There is no total \(D_3\)-equivariant function
\[
 q_\times:\mathcal S_\times\longrightarrow P_2(H_{\rm act})
\]
if \(D_3\) acts on the (26) product signatures through the published coordinate
permutations \(r,s\) and on \(H_{\rm act}\) through
\[
 a(x)=x+3\pmod9,
 \qquad
 j(x)=-x\pmod9.
\]
\end{proposition}

\begin{proof}
The action on \(\mathcal S_\times\) has orbits of sizes
\[
 6+6+6+3+3+1+1.
\]
Its two globally fixed points are the signatures
\(000000\) and \(555555\). The action of \(D_3\) on the six elements of
\(H_{\rm act}\) is regular, and its induced action on the fifteen pairs has
orbits
\[
 6+3+3+3,
\]
with no globally fixed point. An equivariant function must map a fixed point
of the source to a fixed point of the target, which is impossible.
\end{proof}

The obstruction precisely identifies the information lost by the isolated
product signature: the fixed sectors can be resolved only by preserving phase,
orientation, and the enriched record. The preceding encoder \(c_0\) acts on
the two-zero stratum, but it cannot replace the family \(c_k\) with a
single column of \(E_{\rm sig}\).

\section{Visible parity, carry parity and physical boundary}

The even--odd separation can be formulated without depending on a particle
label. For each of the leaves
\(\bullet\in\{+,\times\}\), set
\[
 n_+(i,j)=i+j,
 \qquad
 n_\times(i,j)=ij,
\]
and let us define the two bits
\[
 \epsilon_R^\bullet(i,j)
 =\operatorname{dr}_9(n_\bullet(i,j))\bmod2,
 \qquad
 \epsilon_Q^\bullet(i,j)
 =q_9^\sharp(n_\bullet(i,j))\bmod2.
\]

\begin{theorem}[Parity-Carry Factorization]
\label{thm:factorizacion-paridad-acarreo}
For the \(81\) cells of each APP sheet,
\[
 \boxed{
 n_\bullet(i,j)\bmod2
 =\epsilon_R^\bullet(i,j)+\epsilon_Q^\bullet(i,j)\pmod2.}
\]
The exact census by pairs
\((\epsilon_R,\epsilon_Q)=(00,01,10,11)\) is
\[
 \begin{array}{c|rrrr}
  &00&01&10&11\\
 \hline
 +&16&20&20&25\\
 \times&25&5&20&31
 \end{array}
\]
and, therefore, the additive and multiplicative leaves possess distinct
parity--carry distributions.
\end{theorem}

\begin{proof}
Reduction modulo two
\(n=9q_9^\sharp(n)+\operatorname{dr}_9(n)\) gives the identity, since
\(9\equiv1\pmod2\). On the additive sheet, the multiplicities of \(i+j\) are
\(1,2,\ldots,9,8,\ldots,1\); separating representative and quotient produces the
first row. The same calculation on the \(81\) products \(ij\) produces the
second. Enumeration of the \(162\) cells verifies the identity before
accumulating the counts.\qedhere
\end{proof}

On \(\ell^2(D_9^2)\), the operators diagonal
\[
 \Gamma_R^\bullet f(i,j)
 =(-1)^{\epsilon_R^\bullet(i,j)}f(i,j),
 \qquad
 \Gamma_Q^\bullet f(i,j)
 =(-1)^{\epsilon_Q^\bullet(i,j)}f(i,j)
\]
are commutative involutions. The unreduced parity factorizes as
\[
 \boxed{\Gamma_{\rm raw}^\bullet
 =\Gamma_R^\bullet\Gamma_Q^\bullet.}
\]
This is the precise operatorial formulation of the idea that the digital root
reads the visible component while the quotient preserves the drag of the coronas. The
difference between the two rows of the table is also an exact form of the
incompatibility between addition and product accumulation.

Each involution decomposes the space into even and odd sectors and provides
an internal grading \(\mathbb Z/2\mathbb Z\). A Bose--Einstein or Fermi--Dirac
statistic, however, is an assertion about exchanging several indistinguishable
realizations, not about the parity of a single cell. Its derivation requires
an exchange functor representing symmetric or braid groups
and intertwining the dynamics. HMT parity constructs
the graded domain on which that functor must act; it is not identified
here with its physical representation.
